\documentclass[11pt]{article}
\usepackage[utf8]{inputenc}
\usepackage{amsmath,amssymb}
\usepackage{authblk}
\usepackage[margin=1in]{geometry}
\usepackage[hidelinks]{hyperref}
\usepackage{xurl}
\usepackage{setspace}

\title{Charge Below, Gravity Above:\\[0.2em]
\large RealQM, Maxwell, and Emergent Gravitation}
\author[1]{Claes Johnson}
\affil[1]{KTH Royal Institute of Technology, SE-100 44 Stockholm, Sweden \\
\texttt{claesjohnson@gmail.com}}
\date{\today \\[0.5em]
\small\textit{Written in respectful cooperation between Claes Johnson and Claude (Anthropic), with Claude
contributing a balanced, neutral view.}}

\begin{document}
\maketitle

\begin{abstract}
RealQM together with Maxwell's equations describes matter and light from charge and energy alone: the electron is
a charge cloud in ordinary three-dimensional space, inertia is the derived quantity $E/c^2$, and mass enters
nowhere as a primitive~\cite{NoMass}. Gravitation is absent from this description --- and need not be added. Gravity
is negligible microscopically (between two protons it is $\sim 10^{-36}$ of the electric force) and dominates only
at large scales, where electromagnetism has screened itself away in neutral bulk matter while unscreened gravity
accumulates. We argue that this clean split --- electromagnetic, charged and small \emph{below}; gravitational,
neutral and large \emph{above} --- is naturally read as gravity being \emph{emergent}: a coarse-grained, collective
phenomenon rather than a fundamental force acting between individual particles, in the spirit of Sakharov's induced
gravity~\cite{Sakharov}, Jacobson's derivation of the Einstein equation from thermodynamics~\cite{Jacobson}, and
the holographic emergence of spacetime from entanglement~\cite{VanRaamsdonk,Padmanabhan}. The resulting picture is
economical: \emph{one} mechanics (Newton's second law $\mathbf F=\dot{\mathbf p}$), \emph{one} inertia ($E/c^2$),
and \emph{two} forces separated by regime --- electromagnetism coupling to charge, gravity to energy. On this view inertial mass ($E/c^2$, the energy content)
is fundamental while \emph{gravitational} mass --- the role of that energy as gravity's charge --- is emergent,
which turns the observed equality of the two masses from an unexplained coincidence into a consequence of
emergence. We are explicit about the limits. Emergent gravity emerges as \emph{general relativity}, not Newtonian gravity (Newton is
only its weak-field shadow); the proposal is established only in holographic models, not proven for our universe;
and the framework is confined to non-extreme physics.
\end{abstract}

\setstretch{1.06}

\section{Two forces, two regimes}
Ordinary physics is run by two long-range forces, and they occupy almost disjoint territories. Electromagnetism
governs the small and the charged: atoms, molecules, chemistry, materials, light. Gravitation governs the large
and the neutral: apples, planets, stars, galaxies. The boundary between them is not arbitrary; it follows from two
facts, one on each side.

First, \emph{gravity is extraordinarily weak}. The gravitational attraction between two protons is about
$10^{-36}$ of their electrostatic repulsion; between two electrons, about $10^{-43}$. At the scale of an atom
gravity is not merely small, it is irrelevant to any digit anyone measures. Second, \emph{electromagnetism
screens, and gravity does not}. Charge comes in two signs and bulk matter is neutral, so over macroscopic bodies
the electric forces cancel; mass comes in one sign, so gravity never cancels --- it \emph{accumulates}. Aggregate
enough neutral matter and electromagnetism has quietly removed itself while gravity has been adding up, until, at
planetary scale, the weakest force in nature is the only one left standing.

So the division ``electromagnetic below, gravitational above'' is not a modelling convenience but the way the two
forces are actually organised: weakness keeps gravity out of the microscopic world, and screening keeps
electromagnetism out of the macroscopic one.

\section{One inertia, one mechanics, two forces}
Underneath the two regimes there is a single mechanics. RealQM, together with Maxwell's equations, describes the
microscopic world with \emph{no mass as a primitive}~\cite{NoMass}: the energy of matter is computed from charge
and a geometric scale, and when a charge moves its inertia is the momentum its own field carries,
\begin{equation}
\mathbf p=\frac{E}{c^2}\,\mathbf v ,
\label{eq:inertia}
\end{equation}
so that the inertial ``mass'' is simply the energy content $E/c^2$~\cite{Thomson}. Newton's second law, in its
original form $\mathbf F=\dot{\mathbf p}$, then governs the response, with the force supplied by electromagnetism,
$\mathbf F=q\mathbf E+q\mathbf v\times\mathbf B$, which couples to \emph{charge} and contains no mass at all.

The macroscopic regime uses the \emph{same} law and the \emph{same} inertia. A neutral body has energy $E$, hence
inertia $E/c^2$, and Newton's second law moves it; the only addition is a second force, gravity, which couples not
to charge but to that same energy $E/c^2$. The economy is worth stating plainly:
\begin{center}
\begin{tabular}{lll}
\hline
 & couples to & inertia (response) \\
\hline
electromagnetism & charge $q$ & $E/c^2$ \\
gravitation      & energy $E/c^2$ & $E/c^2$ \\
\hline
\end{tabular}
\end{center}
\emph{One} mechanics ($\mathbf F=\dot{\mathbf p}$), \emph{one} inertia ($E/c^2$), \emph{two} forces reading two
different properties of the same matter --- its charge and its energy. Nothing here needs a separate primitive
called mass; ``mass'' is the name of the inertia, and the inertia is energy.

\section{Why the separation is clean}
Because the two forces read different properties and inhabit different scales, they do not interfere in the
regimes where each operates. In the microscopic world gravity is $10^{-36}$ too weak to enter RealQM's Coulomb
bookkeeping, so RealQM$+$Maxwell is \emph{complete without gravitation}: atoms, bonds, spectra, and the propagation
of light are electromagnetic through and through. In the macroscopic world the constituent charges have screened to
neutrality, so gravitation is \emph{complete without electromagnetism}: celestial mechanics needs no Maxwell field.
The separation is therefore not a truncation we impose but a consequence of \emph{weakness} on one side and
\emph{neutrality} on the other.

There is exactly one regime where the split must fail: where both conditions break at once, i.e.\ where gravity is
strong \emph{and} energy is unscreened. That is the province of astrophysics and cosmology --- light bending near a
star, the field energy of radiation gravitating in the early universe --- and there electromagnetism and
gravitation genuinely couple through their shared currency, energy. Everywhere else, across the whole of
non-extreme physics, the two theories run in disjoint domains.

\section{Gravity as emergent on larger scales}
That gravity should be at once the weakest force, unscreenable, universal in its coupling to energy, and important
only in the large invites a stronger reading than ``fundamental but faint'': that gravity is \emph{emergent} --- a
collective, coarse-grained effect of many degrees of freedom, in the way that temperature and pressure are
collective effects meaningless for a single molecule. This is not a fringe suggestion but a recurring theme of
modern gravitational theory.

Sakharov proposed already in 1967 that gravity is an ``induced'' elastic response of the quantum
vacuum~\cite{Sakharov}. Jacobson showed in 1995 that Einstein's field equations can be \emph{derived} as a
thermodynamic equation of state, from the Clausius relation $\delta Q=T\,dS$ applied to local causal
horizons~\cite{Jacobson} --- gravity appearing not as a fundamental force but as the large-scale thermodynamics of
whatever the underlying degrees of freedom are, much as the gas law follows from statistics without tracking
molecules. Padmanabhan developed this thermodynamic route in detail~\cite{Padmanabhan}, and in the holographic
setting the emergence is sharpest of all: bulk spacetime geometry, and hence gravity, is built from the quantum
\emph{entanglement} of a lower-dimensional field theory~\cite{VanRaamsdonk}.

On this reading the microscopic world has no gravity to speak of --- not merely a negligible amount, but no
fundamental gravitational interaction between individual charges at all, just as a single molecule has no
temperature. Gravity is what \emph{emerges} when enough energy is coarse-grained together. This fits the two-regime
picture exactly: the fundamental layer is charge and energy (RealQM$+$Maxwell), and gravitation is the large-scale
collective phenomenon that appears above it. The very features that made the separation clean in \S3 --- gravity's
weakness, its non-screening, its universal coupling to energy --- are precisely what one expects of an emergent,
statistical effect rather than a microscopic force.

\section{Gravitational mass: an emergent role of a fundamental quantity}
A corollary sharpens what is, and is not, emergent, and it concerns mass directly. The \emph{quantity} $E/c^2$ ---
the energy content of a piece of matter --- is fundamental and present at every scale: it is the inertial mass,
RealQM computes it for a single atom, and a lone electron carries it. What is emergent is the \emph{role} of that
quantity as \emph{gravitational} mass, the charge that sources and feels gravity. Where there is no gravity ---
the microscopic world --- $E/c^2$ is inertia and nothing more; gravitational mass is not merely negligible there,
it has \emph{no referent}, exactly as a single molecule has no temperature. Gravitational mass appears only when
gravity does, at large scales. So inertial mass is fundamental; gravitational mass is the emergent role it takes on
once gravity has emerged.

This has a welcome consequence: it \emph{explains} the equivalence principle rather than positing it. The equality
of gravitational and inertial mass --- the deepest empirical fact about gravity, tested to $10^{-13}$ --- is
ordinarily an unexplained coincidence. Here it is automatic. Emergent gravity arises as the collective response to
the energy that is \emph{already present} as inertia, so it can couple to nothing but $E/c^2$: the gravitational
charge \emph{is} the inertial energy, by construction and not by accident. Equivalence is then not an input but an
output of emergence.

\section{What is honest to claim, and what is not}
The picture is attractive and, in its non-extreme domain, coherent; but three limits must be stated plainly rather
than glossed.

\emph{First, emergent gravity emerges as general relativity, not Newton.} Jacobson's construction yields Einstein's
equations, and holography yields general relativity in the bulk. A consistent large-scale gravity that couples to
energy and propagates at finite speed is, in any case, forced towards general relativity by the requirement that
the gravitational field's own energy also gravitate. Newtonian gravity is only the weak-field, slow-motion corner
of what emerges; it is the working form in \S2, but not the fundamental large-scale theory, which must reproduce
light bending, perihelion precession and gravitational waves.

\emph{Second, emergence is proven only in models, not for our world.} The thermodynamic and holographic derivations
are compelling and, in anti-de~Sitter/conformal-field-theory examples, essentially established; but that gravity is
emergent \emph{in our universe} remains a hypothesis. The standard framework still treats gravity as a fundamental
(spin-2) field, and the two views are not yet experimentally separated.

\emph{Third, the framework is bounded to non-extreme physics, and inherits open sectors.} It says nothing about the
strong and weak nuclear forces except through RealQM's own (unconventional) charge-cloud modelling of nuclei; the
elementary rest energies remain inputs, as in every theory; and the genuinely relativistic and quantum-field
regimes --- pair creation, the precision anomalies of quantum electrodynamics, and quantum gravity itself --- lie
outside it by construction. Naming that boundary is a strength: it excludes exactly the regimes where the clean
separation is known to break.

\section{Conclusion}
For the vast, non-extreme middle of physics --- chemistry, materials, optics, mechanics, the motion of ordinary
bodies --- a single economical picture suffices. The fundamental layer is \emph{charge and energy}: RealQM and
Maxwell, with inertia the derived quantity $E/c^2$ and no primitive mass. Upon it act two forces reading two
properties of the same matter: electromagnetism its charge, gravitation its energy, kept in disjoint domains by the
weakness of gravity below and the neutrality of bulk matter above. And the most natural place for gravitation in
this picture is not as a second fundamental force but as an \emph{emergent}, large-scale, collective effect ---
appearing, like thermodynamics, only when enough energy is gathered together, and absent as a microscopic
interaction between charges. What emerges must, to match the sky, be general relativity rather than Newton, and
whether it truly emerges in our universe is still open; but the organising thought is clean and, within its
declared bounds, self-consistent: \emph{charge and energy below, emergent gravity above, and a single inertia
$E/c^2$ carried through both.}

\begin{thebibliography}{9}
\bibitem{NoMass} C.~Johnson, \emph{Energy Without Mass: RealQM, Inertia, and the Penning Trap} (companion note).
\bibitem{RealQM} C.~Johnson, \emph{Real Quantum Mechanics} (book).
\bibitem{Thomson} J.~J.~Thomson, ``On the electric and magnetic effects produced by the motion of electrified
  bodies,'' \emph{Philos.\ Mag.} \textbf{11}, 229--249 (1881).
\bibitem{Sakharov} A.~D.~Sakharov, ``Vacuum quantum fluctuations in curved space and the theory of gravitation,''
  \emph{Sov.\ Phys.\ Dokl.} \textbf{12}, 1040--1041 (1968).
\bibitem{Jacobson} T.~Jacobson, ``Thermodynamics of spacetime: the Einstein equation of state,''
  \emph{Phys.\ Rev.\ Lett.} \textbf{75}, 1260--1263 (1995).
\bibitem{Padmanabhan} T.~Padmanabhan, ``Thermodynamical aspects of gravity: new insights,''
  \emph{Rep.\ Prog.\ Phys.} \textbf{73}, 046901 (2010).
\bibitem{VanRaamsdonk} M.~Van~Raamsdonk, ``Building up spacetime with quantum entanglement,''
  \emph{Gen.\ Rel.\ Grav.} \textbf{42}, 2323--2329 (2010).
\bibitem{MMR} C.~Johnson, \emph{Many-Minds Relativity} (book).
\end{thebibliography}

\end{document}
